\begin{center}
\LARGE
Computer Algebra application to the distribution of sample
                     correlation coefficient

\end{center}

\begin{center}
\Large
Shigekazu Nakagawa, Naoto Niki, Hiroki Hashiguchi
\end{center}

\noindent
Date:  July 18th (Thursday)\\
Time:  16:00-16:25\\

\begin{center}
\large
Abstract
\end{center}

Let $r$ be the correlation coefficient based on a sample of size $n$ from a population $F$. The
value of $r$ remains unchanged however observations may be exchanged or permuted. The
statistic $r$ is one of symmetric statistics. 

Our goal is divided into two parts: 

\begin{enumerate}

\item
obtaining the asymptotic cumulants of $r$,
\item
as a consequence of 1., deriving the asymptotic expansions for probability integrals and
    percentiles of $r$, where $F$ has population cumulants of requisite order. 


\end{enumerate}


In accomplishing 1, the computational difficulty arises there. If we want the more precise
approximation, the order of requisite approximate cumulants of $r$ increases. Raising the order
may cost a large amount of algebraic computation often beyond human power. Here computer
algebra induces an increasing interest. 

The authors have taken a new look at symmetric statistics and developed the symbolic and
algebraic algorithm for changing of bases of symmetric polynomials. Both the asymptotic
cumulants and the approximate distributions are given. 



